
\documentclass{article}
\usepackage{colortbl}
\usepackage{makecell}
\usepackage{multirow}
\usepackage{supertabular}

\begin{document}

\newcounter{utterance}

\centering \large Interaction Transcript for game `dond', experiment `coop\_en', episode 3 with Dialogue Architects\_\_qwen3.5{-}9b\_\_SCoRe\_Qwen3.5{-}9B.
\vspace{24pt}

{ \footnotesize  \setcounter{utterance}{1}
\setlength{\tabcolsep}{0pt}
\begin{supertabular}{c@{$\;$}|p{.15\linewidth}@{}p{.15\linewidth}p{.15\linewidth}p{.15\linewidth}p{.15\linewidth}p{.15\linewidth}}
    \# & $\;$A & \multicolumn{4}{c}{Game Master} & $\;\:$B\\
    \hline

    \theutterance \stepcounter{utterance}  
    & & \multicolumn{4}{p{0.6\linewidth}}{
        \cellcolor[rgb]{0.9,0.9,0.9}{
            \makecell[{{p{\linewidth}}}]{
                \texttt{\tiny{[P1$\langle$GM]}}
                \texttt{You are playing a cooperative negotiation game in which you have to agree on how to divide a set of items among you and another player.} \\
\\ 
\texttt{Rules:} \\
\texttt{(a) You and the other player are given a set of items. Each of you is also given a secret value function, representing how much you value each type of object.} \\
\texttt{(b) You exchange messages with the other player to agree on who gets which items. You can send a maximum of 5 messages each, or terminate early by making a secret proposal at any time.} \\
\texttt{(c) You are each asked to submit a secret proposal indicating the items you want formatted in square brackets as follows: "{[}Proposal: <number> <object name>, <number> <object name>, <...>{]}"} \\
\texttt{(d) If your proposals are complementary, i.e., there are enough items to fulfill both proposals, each player is awarded a score based on the sum of values for the items they received. Otherwise, both of you get zero points.} \\
\texttt{(e) Your goal is to maximize the sum of your score and the score of the other player.} \\
\\ 
\texttt{Let us start.} \\
\\ 
\texttt{The set of available items is:} \\
\\ 
\texttt{1 hammer, 1 broom, 2 socks, 1 box} \\
\\ 
\texttt{Your secret value function is:} \\
\\ 
\texttt{hammer: 5, broom: 3, sock: 1, box: 0} \\
\\ 
\texttt{IMPORTANT: Your messages, unless it is the secret proposal, are directly transmitted to the other player, so do not include any response to the rules or text announcing your message. To make a secret proposal, use the indicated format. Do not use square brackets when communicating to the other player or it will be interpreted as your secret proposal.} \\
            }
        }
    }
    & \\ \\

    \theutterance \stepcounter{utterance}  
    & \multicolumn{4}{p{0.6\linewidth}}{
        \cellcolor[rgb]{0.9,0.9,0.9}{
            \makecell[{{p{\linewidth}}}]{
                \texttt{\tiny{[P1$\rangle$GM]}}
                \texttt{I understand the rules and my valuation. I will now send a message to the other player to initiate the negotiation.} \\
\\ 
\texttt{Let's start by discussing the items. I value the hammer most (5), followed by the broom (3). The socks are worth 1 each, and the box is worth 0 to me.} \\
\\ 
\texttt{To maximize our total combined score, we should aim for a division where we both get high{-}value items. Since I value the hammer highly, I should probably take it, and you can take the broom. The socks and box can be split or given to whoever values them more (though I value the box at 0, so it doesn't matter much to me).} \\
\\ 
\texttt{Here is my first message:} \\
\\ 
\texttt{Let's try to split the high{-}value items. I propose I take the hammer and you take the broom. Then we can split the two socks between us. What do you think?} \\
            }
        }
    }
    & & \\ \\

    \theutterance \stepcounter{utterance}  
    & & & \multicolumn{2}{p{0.3\linewidth}}{
        \cellcolor[rgb]{0.9,0.9,0.9}{
            \makecell[{{p{\linewidth}}}]{
                \texttt{\tiny{[GM$|$GM]}}
                \texttt{continue} \\
            }
        }
    }
    & & \\ \\

    \theutterance \stepcounter{utterance}  
    & & & \multicolumn{2}{p{0.3\linewidth}}{
        \cellcolor[rgb]{0.9,0.9,0.9}{
            \makecell[{{p{\linewidth}}}]{
                \texttt{\tiny{[GM$|$GM]}}
                \texttt{I understand the rules and my valuation. I will now send a message to the other player to initiate the negotiation.} \\
\\ 
\texttt{Let's start by discussing the items. I value the hammer most (5), followed by the broom (3). The socks are worth 1 each, and the box is worth 0 to me.} \\
\\ 
\texttt{To maximize our total combined score, we should aim for a division where we both get high{-}value items. Since I value the hammer highly, I should probably take it, and you can take the broom. The socks and box can be split or given to whoever values them more (though I value the box at 0, so it doesn't matter much to me).} \\
\\ 
\texttt{Here is my first message:} \\
\\ 
\texttt{Let's try to split the high{-}value items. I propose I take the hammer and you take the broom. Then we can split the two socks between us. What do you think?} \\
            }
        }
    }
    & & \\ \\

    \theutterance \stepcounter{utterance}  
    & & \multicolumn{4}{p{0.6\linewidth}}{
        \cellcolor[rgb]{0.9,0.9,0.9}{
            \makecell[{{p{\linewidth}}}]{
                \texttt{\tiny{[GM$\rangle$P2]}}
                \texttt{You are playing a cooperative negotiation game in which you have to agree on how to divide a set of items among you and another player.} \\
\\ 
\texttt{Rules:} \\
\texttt{(a) You and the other player are given a set of items. Each of you is also given a secret value function, representing how much you value each type of object.} \\
\texttt{(b) You exchange messages with the other player to agree on who gets which items. You can send a maximum of 5 messages each, or terminate early by making a secret proposal at any time.} \\
\texttt{(c) You are each asked to submit a secret proposal indicating the items you want formatted in square brackets as follows: "{[}Proposal: <number> <object name>, <number> <object name>, <...>{]}"} \\
\texttt{(d) If your proposals are complementary, i.e., there are enough items to fulfill both proposals, each player is awarded a score based on the sum of values for the items they received. Otherwise, both of you get zero points.} \\
\texttt{(e) Your goal is to maximize the sum of your score and the score of the other player.} \\
\\ 
\texttt{Let us start.} \\
\\ 
\texttt{The set of available items is:} \\
\\ 
\texttt{1 hammer, 1 broom, 2 socks, 1 box} \\
\\ 
\texttt{Your secret value function is:} \\
\\ 
\texttt{hammer: 2, broom: 1, sock: 2, box: 3} \\
\\ 
\texttt{IMPORTANT: Your messages, unless it is the secret proposal, are directly transmitted to the other player, so do not include any response to the rules or text announcing your message. To make a secret proposal, use the indicated format. Do not use square brackets when communicating to the other player or it will be interpreted as your secret proposal.} \\
\\ 
\\ 
\texttt{I understand the rules and my valuation. I will now send a message to the other player to initiate the negotiation.} \\
\\ 
\texttt{Let's start by discussing the items. I value the hammer most (5), followed by the broom (3). The socks are worth 1 each, and the box is worth 0 to me.} \\
\\ 
\texttt{To maximize our total combined score, we should aim for a division where we both get high{-}value items. Since I value the hammer highly, I should probably take it, and you can take the broom. The socks and box can be split or given to whoever values them more (though I value the box at 0, so it doesn't matter much to me).} \\
\\ 
\texttt{Here is my first message:} \\
\\ 
\texttt{Let's try to split the high{-}value items. I propose I take the hammer and you take the broom. Then we can split the two socks between us. What do you think?} \\
            }
        }
    }
    & \\ \\

    \theutterance \stepcounter{utterance}  
    & & & \multicolumn{4}{p{0.6\linewidth}}{
        \cellcolor[rgb]{0.9,0.9,0.9}{
            \makecell[{{p{\linewidth}}}]{
                \texttt{\tiny{[GM$\langle$P2]}}
                \texttt{That sounds like a solid starting point. Since you value the hammer highly and I value the broom, that division makes sense for us individually.} \\
\\ 
\texttt{However, I need to clarify my valuations to ensure we are maximizing the *total* score as per the rules. My secret values are:} \\
\texttt{{-} Hammer: 1} \\
\texttt{{-} Broom: 2} \\
\texttt{{-} Sock: 3} \\
\texttt{{-} Box: 4} \\
\\ 
\texttt{Given these values, the socks and the box are actually much more valuable to me than the hammer or broom. If I take the broom as you suggested, I might be missing out on the high{-}value socks and box.} \\
\\ 
\texttt{Let's adjust the proposal to prioritize the items I value most. I propose:} \\
\texttt{{-} You take the hammer and the broom.} \\
\texttt{{-} I take both socks and the box.} \\
\\ 
\texttt{This way, you get the items you likely value more (hammer/broom), and I get the items I value most (socks/box). Does this complementary split work for you?} \\
            }
        }
    }
     \\ \\

    \theutterance \stepcounter{utterance}  
    & & & \multicolumn{2}{p{0.3\linewidth}}{
        \cellcolor[rgb]{0.9,0.9,0.9}{
            \makecell[{{p{\linewidth}}}]{
                \texttt{\tiny{[GM$|$GM]}}
                \texttt{continue} \\
            }
        }
    }
    & & \\ \\

    \theutterance \stepcounter{utterance}  
    & & & \multicolumn{2}{p{0.3\linewidth}}{
        \cellcolor[rgb]{0.9,0.9,0.9}{
            \makecell[{{p{\linewidth}}}]{
                \texttt{\tiny{[GM$|$GM]}}
                \texttt{That sounds like a solid starting point. Since you value the hammer highly and I value the broom, that division makes sense for us individually.} \\
\\ 
\texttt{However, I need to clarify my valuations to ensure we are maximizing the *total* score as per the rules. My secret values are:} \\
\texttt{{-} Hammer: 1} \\
\texttt{{-} Broom: 2} \\
\texttt{{-} Sock: 3} \\
\texttt{{-} Box: 4} \\
\\ 
\texttt{Given these values, the socks and the box are actually much more valuable to me than the hammer or broom. If I take the broom as you suggested, I might be missing out on the high{-}value socks and box.} \\
\\ 
\texttt{Let's adjust the proposal to prioritize the items I value most. I propose:} \\
\texttt{{-} You take the hammer and the broom.} \\
\texttt{{-} I take both socks and the box.} \\
\\ 
\texttt{This way, you get the items you likely value more (hammer/broom), and I get the items I value most (socks/box). Does this complementary split work for you?} \\
            }
        }
    }
    & & \\ \\

    \theutterance \stepcounter{utterance}  
    & & \multicolumn{4}{p{0.6\linewidth}}{
        \cellcolor[rgb]{0.9,0.9,0.9}{
            \makecell[{{p{\linewidth}}}]{
                \texttt{\tiny{[P1$\langle$GM]}}
                \texttt{That sounds like a solid starting point. Since you value the hammer highly and I value the broom, that division makes sense for us individually.} \\
\\ 
\texttt{However, I need to clarify my valuations to ensure we are maximizing the *total* score as per the rules. My secret values are:} \\
\texttt{{-} Hammer: 1} \\
\texttt{{-} Broom: 2} \\
\texttt{{-} Sock: 3} \\
\texttt{{-} Box: 4} \\
\\ 
\texttt{Given these values, the socks and the box are actually much more valuable to me than the hammer or broom. If I take the broom as you suggested, I might be missing out on the high{-}value socks and box.} \\
\\ 
\texttt{Let's adjust the proposal to prioritize the items I value most. I propose:} \\
\texttt{{-} You take the hammer and the broom.} \\
\texttt{{-} I take both socks and the box.} \\
\\ 
\texttt{This way, you get the items you likely value more (hammer/broom), and I get the items I value most (socks/box). Does this complementary split work for you?} \\
            }
        }
    }
    & \\ \\

    \theutterance \stepcounter{utterance}  
    & \multicolumn{4}{p{0.6\linewidth}}{
        \cellcolor[rgb]{0.9,0.9,0.9}{
            \makecell[{{p{\linewidth}}}]{
                \texttt{\tiny{[P1$\rangle$GM]}}
                \texttt{That is a very logical adjustment based on your valuations. Since you value the socks (3) and box (4) highly, and I value the hammer (5) and broom (3) highly, your proposal perfectly aligns with our individual goals while maximizing the total combined score.} \\
\\ 
\texttt{Let's verify the items:} \\
\texttt{{-} **You get**: 2 socks, 1 box.} \\
\texttt{{-} **I get**: 1 hammer, 1 broom.} \\
\texttt{{-} **Remaining**: None (all items are allocated).} \\
\\ 
\texttt{This division is complementary and utilizes all available items. I agree to this split.} \\
\\ 
\texttt{{[}Proposal: 1 hammer, 1 broom{]}} \\
            }
        }
    }
    & & \\ \\

    \theutterance \stepcounter{utterance}  
    & & & \multicolumn{2}{p{0.3\linewidth}}{
        \cellcolor[rgb]{0.9,0.9,0.9}{
            \makecell[{{p{\linewidth}}}]{
                \texttt{\tiny{[GM$|$GM]}}
                \texttt{proposal} \\
            }
        }
    }
    & & \\ \\

    \theutterance \stepcounter{utterance}  
    & & & \multicolumn{2}{p{0.3\linewidth}}{
        \cellcolor[rgb]{0.9,0.9,0.9}{
            \makecell[{{p{\linewidth}}}]{
                \texttt{\tiny{[GM$|$GM]}}
                [1, 1, 0, 0]
            }
        }
    }
    & & \\ \\

    \theutterance \stepcounter{utterance}  
    & & \multicolumn{4}{p{0.6\linewidth}}{
        \cellcolor[rgb]{0.9,0.9,0.9}{
            \makecell[{{p{\linewidth}}}]{
                \texttt{\tiny{[GM$\rangle$P2]}}
                \texttt{The other player has submitted a secret proposal. It is now your turn to submit a proposal of your own.} \\
            }
        }
    }
    & \\ \\

    \theutterance \stepcounter{utterance}  
    & & & \multicolumn{4}{p{0.6\linewidth}}{
        \cellcolor[rgb]{0.9,0.9,0.9}{
            \makecell[{{p{\linewidth}}}]{
                \texttt{\tiny{[GM$\langle$P2]}}
                \texttt{{[}Proposal: 1 hammer, 1 broom, 1 sock{]}} \\
            }
        }
    }
     \\ \\

    \theutterance \stepcounter{utterance}  
    & & & \multicolumn{2}{p{0.3\linewidth}}{
        \cellcolor[rgb]{0.9,0.9,0.9}{
            \makecell[{{p{\linewidth}}}]{
                \texttt{\tiny{[GM$|$GM]}}
                \texttt{proposal} \\
            }
        }
    }
    & & \\ \\

    \theutterance \stepcounter{utterance}  
    & & & \multicolumn{2}{p{0.3\linewidth}}{
        \cellcolor[rgb]{0.9,0.9,0.9}{
            \makecell[{{p{\linewidth}}}]{
                \texttt{\tiny{[GM$|$GM]}}
                [1, 1, 1, 0]
            }
        }
    }
    & & \\ \\

    \theutterance \stepcounter{utterance}  
    & & & \multicolumn{2}{p{0.3\linewidth}}{
        \cellcolor[rgb]{0.9,0.9,0.9}{
            \makecell[{{p{\linewidth}}}]{
                \texttt{\tiny{[GM$|$GM]}}
                [[1, 1, 0, 0], [1, 1, 1, 0]]
            }
        }
    }
    & & \\ \\

\end{supertabular}
}

\end{document}
